\subsection{The quantum route, and the window as the bounds draw it}
\label{sec:quantum}

The Hamiltonian~\eqref{eq:HN} on the grid is a sparse, local lattice model, the class that quantum computers simulate with polynomial overhead~\cite{lloyd1996,zalka1998,kassal2008}. We price it with qubitization~\cite{low2019qubitization} and, for comparison, a product formula~\cite{childs2021trotter} at precision $\varepsilon=\num{eps_prec}$ under the budgets of Table~\ref{tab:params}. The counts are leading-order formulas (Appendix~\ref{app:quantum}), not compiled circuits, and they price \emph{one} coherent evolution to $\Tobs$, without state preparation, repeated measurements or error-correction overheads beyond the T-gate count. The register grows linearly in $N$ and the qubit budget allows $N_Q=\num{NQ}$; at the simulated mass qubitization stays far inside the T-gate budget for every $N$ of Table~\ref{tab:quantum} \stM{}, while the product formula costs orders of magnitude more and exceeds it already at $N=4$.

\begin{table}[ht]
\centering
\footnotesize
\caption{Quantum price as a function of cluster size (leading-order model, Appendix~\ref{app:quantum}). T-gates at the simulated mass $\mut=\num{mut0}$; $\mutq$: largest mass whose qubitization count fits the T-gate budget (leading order, and with the two omitted terms restored).}
\label{tab:quantum}
\begin{tabular}{rrrrrr}
\toprule
$N$ & qubits & T-gates, qubitization & T-gates, Trotter & $\mutq$ & $\mutq$, complete \\
\midrule
4 & 74 & \ensuremath{1.3\times10^{8}} & \ensuremath{1.6\times10^{12}} & 22.2 & 20.5 \\
8 & 118 & \ensuremath{3.8\times10^{8}} & \ensuremath{1.6\times10^{13}} & 20.5 & 19.0 \\
16 & 206 & \ensuremath{1.3\times10^{9}} & \ensuremath{1.4\times10^{14}} & 18.6 & 17.3 \\
24 & 294 & \ensuremath{2.7\times10^{9}} & \ensuremath{4.9\times10^{14}} & 17.4 & 16.2 \\
40 & 470 & \ensuremath{7.2\times10^{9}} & \ensuremath{2.3\times10^{15}} & 15.9 & 14.8 \\
88 & 998 & \ensuremath{3.3\times10^{10}} & \ensuremath{2.5\times10^{16}} & 13.5 & 12.7 \\
\bottomrule
\end{tabular}
\end{table}


The question is therefore how heavy the particles can get. The largest mass whose qubitization count fits the budget is $\mutq(N)$ (Eq.~\eqref{eq:muqub}; restoring two omitted terms, Appendix~\ref{app:quantum}, lowers it somewhat, Table~\ref{tab:quantum}). It moves only logarithmically with the budget (a thousandfold more T gates gains $\num{mu_budget_shift}$ in $\mut$) and falls slowly with $N$, and at the edge the price is movie length, not qubits. In an optical lattice $\mut=3\sqrt s$ for every species ($s$ in recoil energies), so the single-particle edge at $N=8$ corresponds to about $\num{s_edge_8}$ recoil energies; the pair cost differs between species and is not priced.

Fig.~\ref{fig:regime} draws the two bounds together: the window runs from the tensor-network wall (a band, Section~\ref{sec:walls}) up to $N_Q$ at the simulated mass, and its upper end moves in at larger mass. The quantum edge is the better-defined of the two. Neither is a statement about what the task requires: the price is for one evolution, not a hardness result, and the classical costs are those of the generic routes, not of a method that exploits the two-level structure of Section~\ref{sec:static} or of imaginary-time Monte Carlo (Section~\ref{sec:mc}). If the conditional reduction (Proposition~\ref{prop:cond}) holds with small $M$ at large $N$, the cost that matters is that of the static spectrum, which Table~\ref{tab:quantum} does not price.

